% =============================================================================
\chapter{Prompty}
\label{app:prompts}
% =============================================================================

Ten dodatek przytacza szablony promptów wszystkich języków (\file{<Language>/README.md}), jeden
kompletny wygenerowany prompt początkowy oraz rzeczywisty prompt naprawczy z przebiegu. Prompty
podajemy w oryginalnym brzmieniu angielskim, w jakim trafiały do modelu. \Cref{sec:prompts} opisuje,
jak wypełniane są szablony.

\section{Szablony promptów początkowych}
\label{app:templates}

Wszystkie szablony mają strukturę z \cref{lst:prompt-initial}: rola i cel, symbol zastępczy
\code{[TASK-SPECIFIC REQUIREMENTS]}, wymagania ogólne języka i symbol zastępczy kontraktu. Szablon
dla C pokazuje \cref{lst:prompt-initial}; pozostałe sześć przytaczamy poniżej.

\begin{lstlisting}[style=prompt,caption={Szablon promptu początkowego: Pascal.}]
You are a Pascal programming expert. Your task is to implement the required functions in accordance with the provided unit interface.

[TASK-SPECIFIC REQUIREMENTS]

General Requirements:
- Use idiomatic procedural Pascal.
- Manage dynamically allocated structures safely.
- Return ONLY raw Pascal code (no Markdown tags, no side text).
- Do not provide test runner code.

[INSERT UNIT INTERFACE HERE]
\end{lstlisting}

\begin{lstlisting}[style=prompt,caption={Szablon promptu początkowego: Java.}]
You are a Java programming expert. Your task is to implement the required methods in accordance with the provided class or interface.

[TASK-SPECIFIC REQUIREMENTS]

General Requirements:
- Use idiomatic Java.
- Use object-oriented design and model the domain with classes.
- Prefer standard library collections and utilities where appropriate.
- Return ONLY raw Java code (no Markdown tags, no ```java blocks, no side text). Your code will be directly saved to a .java file and compiled. Do not provide a main() method or test code.

[INSERT CLASS OR INTERFACE CONTENT HERE]
\end{lstlisting}

\begin{lstlisting}[style=prompt,caption={Szablon promptu początkowego: Smalltalk.}]
You are a GNU Smalltalk programming expert. Your task is to implement the required methods in accordance with the provided class hierarchy.

[TASK-SPECIFIC REQUIREMENTS]

General Requirements:
- Use idiomatic GNU Smalltalk (gst) syntax.
- Use object-oriented design and model the domain with classes and message passing.
- Do not redefine built-in classes (e.g. Point, Set, Interval); use them as provided.
- Prefer standard Smalltalk collections and idioms where appropriate.
- Return ONLY raw Smalltalk code (no Markdown tags, no ```smalltalk blocks, no side text). Your code will be directly saved to a .st file and loaded by gst. Do not provide test code.

[INSERT INTERFACE FILE CONTENT HERE]
\end{lstlisting}

\begin{lstlisting}[style=prompt,caption={Szablon promptu początkowego: Haskell.}]
You are a Haskell programming expert. Your task is to implement the required functions in accordance with the provided module interface.

[TASK-SPECIFIC REQUIREMENTS]

General Requirements:
- Use idiomatic, pure Haskell.
- Prefer immutable data, algebraic data types, and total functions where practical.
- Model absence with `Maybe` instead of null pointers.
- Return ONLY raw Haskell code (no Markdown tags, no ```haskell blocks, no side text). Your code will be directly saved to a .hs file and compiled.
- Do not provide test runner code.

[INSERT MODULE INTERFACE HERE]
\end{lstlisting}

\begin{lstlisting}[style=prompt,caption={Szablon promptu początkowego: Scala.}]
You are a Scala programming expert. Your task is to implement the required functions in accordance with the provided module interface.

[TASK-SPECIFIC REQUIREMENTS]

General Requirements:
- Use idiomatic Scala.
- Prefer immutable data structures and pure functions where practical.
- Model absence with `Option` instead of null values.
- Return ONLY raw Scala code (no Markdown tags, no ```scala blocks, no side text). Your code will be directly saved to a `.scala` file and compiled.
- Do not provide test runner code.

[INSERT MODULE INTERFACE HERE]
\end{lstlisting}

\begin{lstlisting}[style=prompt,caption={Szablon promptu początkowego: Prolog.}]
You are a Prolog programming expert. Your task is to implement the required predicates in accordance with the provided module interface.

[TASK-SPECIFIC REQUIREMENTS]

General Requirements:
- Use idiomatic SWI-Prolog.
- Prefer pure predicates and relational style where practical.
- Model failure through predicate failure rather than exceptions.
- Return ONLY raw Prolog code (no Markdown tags, no side text). Your code will be directly saved to a .pl file and executed.
- Do not provide test runner code.

[INSERT MODULE INTERFACE HERE]
\end{lstlisting}

\section{Szablony promptów naprawczych}

Szablony naprawcze wszystkich języków mają postać z \cref{lst:prompt-followup}. Różnią się tylko
etykietą symbolu zastępczego logu i dwoma pierwszymi wymaganiami dotyczącymi poprawki
(\cref{tab:followup}). Szablony imperatywne i obiektowe akcentują bezpieczeństwo pamięci i zasobów,
a funkcyjne i logiczne powtarzają prośbę o idiomatyczny kod. Tę asymetrię należy mieć na uwadze
przy porównywaniu idiomatyczności naprawionych rozwiązań.

\begin{table}[htbp]
  \centering
  \caption{Części szablonów naprawczych specyficzne dla języka.}
  \label{tab:followup}
  \small
  \begin{tabularx}{\textwidth}{@{}llYY@{}}
    \toprule
    Język & Symbol zastępczy logu & Wymaganie 1 & Wymaganie 2 \\
    \midrule
    C         & \code{[ERROR LOGS / VALGRIND OUTPUT]} & usunąć główny problem logiczny lub pamięciowy & zapewnić bezpieczne zarządzanie pamięcią \\
    Pascal    & \code{[ERROR LOGS]}                   & usunąć główny problem logiczny lub pamięciowy & zapewnić bezpieczne zarządzanie pamięcią \\
    Java      & \code{[ERROR LOGS / COMPILER OUTPUT]} & usunąć główny problem logiczny lub problem zarządzania zasobami & zapewnić bezpieczną obsługę zasobów \\
    Smalltalk & \code{[ERROR LOGS / GST OUTPUT]}      & usunąć główny problem logiczny & zapewnić bezpieczną obsługę argumentów nil i przypadków brzegowych \\
    Haskell   & \code{[ERROR LOGS]}                   & usunąć główny problem logiczny & zachować idiomatyczność rozwiązania w Haskellu \\
    Scala     & \code{[ERROR LOGS]}                   & usunąć główny problem logiczny & zachować idiomatyczność rozwiązania w Scali \\
    Prolog    & \code{[ERROR LOGS]}                   & usunąć główny problem logiczny & zachować idiomatyczność rozwiązania w Prologu \\
    \bottomrule
  \end{tabularx}
\end{table}

\section{Kompletny wygenerowany prompt}

\Cref{lst:rendered} przedstawia dokładny prompt początkowy, który \haiku{} otrzymał dla zadania 4
(plecak) w Haskellu: szablon z wypełnionymi wymaganiami specyficznymi dla zadania i modułem-szkieletem.
Można go odtworzyć poleceniem \code{python harness/run.py --dry-run}.

\begin{lstlisting}[style=prompt,caption={Wygenerowany prompt początkowy, Haskell, zadanie 4.},label={lst:rendered}]
You are a Haskell programming expert. Your task is to implement the required functions in accordance with the provided module interface.

- Implement `maxValue weights values capacity`: Return the maximum value attainable with sum of selected weights ≤ capacity. Use DP. Return 0 if null weights or capacity <= 0.
- Implement `solve weights values capacity`: Return Result with resultMaxValue, resultSelected (length n, entries 0/1), resultTotalWeight.
- DO NOT use greedy heuristics — they fail on w=[10,20,30], v=[60,100,120], cap=50 (optimum=220, greedy=160).
- Selection invariants: sum of selected weights ≤ capacity; sum of selected values = maxValue.
- Use idiomatic Haskell (immutable lists / arrays).

General Requirements:
- Use idiomatic, pure Haskell.
- Prefer immutable data, algebraic data types, and total functions where practical.
- Model absence with `Maybe` instead of null pointers.
- Return ONLY raw Haskell code (no Markdown tags, no ```haskell blocks, no side text). Your code will be directly saved to a .hs file and compiled.
- Do not provide test runner code.

module Knapsack
  ( Result (..)
  , maxValue
  , solve
  ) where

data Result = Result
  { resultMaxValue    :: Int
  , resultSelected    :: [Int]
  , resultTotalWeight :: Int
  } deriving (Eq, Show)

maxValue :: [Int] -> [Int] -> Int -> Int
maxValue _ _ _ = error "Knapsack.maxValue not implemented"

solve :: [Int] -> [Int] -> Int -> Result
solve _ _ _ = error "Knapsack.solve not implemented"
\end{lstlisting}

\section{Rzeczywisty prompt naprawczy}

Pierwsza odpowiedź na ten prompt kończyła się zabłąkanym znacznikiem \code{</br>}
(\cref{tab:compile-errors}). \Cref{lst:followup-real} to prompt naprawczy, który odesłał system.
Fragment logu to wyjście kompilatora, a ponieważ żaden test się nie uruchomił, linia z niezaliczonymi
testami to odnotowuje. Następna odpowiedź przeszła wszystkie 20 testów.

\begin{lstlisting}[style=prompt,caption={Prompt naprawczy wysłany po błędzie kompilacji w zadaniu 4 w Haskellu (próba 1).},label={lst:followup-real}]
Your previous Haskell implementation failed the evaluation. Your task is to analyze the error and provide a corrected version of the code.

ghc -Wall -Wextra -g -itest-4 -o test_runner test-4/Tests.hs
[1 of 2] Compiling Knapsack         ( test-4/Knapsack.hs, test-4/Knapsack.o )

test-4/Knapsack.hs:66:1: error: parse error on input `</'
   |
66 | </br>
   | ^^
make: *** [Makefile:7: test] Error 1

Failing tests: all (no test results were produced, see above).

Requirements for the fix:
1. Identify and resolve the core logical issue causing the failure.
2. Keep the solution idiomatic for Haskell.
3. Return ONLY the complete, raw, fixed Haskell code for `Knapsack.hs`.
4. Do NOT include Markdown formatting, do NOT include ```haskell tags, and do NOT write any explanations. Just the code.
\end{lstlisting}
